\documentclass[12pt,a4paper]{article}
\usepackage[utf8]{inputenc}
\usepackage[T1]{fontenc}
\usepackage[french]{babel}
\usepackage[margin=2.5cm]{geometry}
\usepackage{xcolor}
\usepackage{listings}
\usepackage{hyperref}
\usepackage{booktabs}
\usepackage{tabularx}
\usepackage{longtable}
\usepackage{array}
\usepackage{fancyhdr}
\usepackage{titlesec}
\usepackage{mdframed}
\usepackage{enumitem}
\usepackage{amsmath}

% ---- Couleurs ----
\definecolor{primary}{RGB}{26, 86, 219}
\definecolor{secondary}{RGB}{30, 64, 175}
\definecolor{success}{RGB}{5, 150, 105}
\definecolor{danger}{RGB}{220, 38, 38}
\definecolor{codebg}{RGB}{245, 247, 250}
\definecolor{codeborder}{RGB}{203, 213, 225}
\definecolor{codecomment}{RGB}{107, 114, 128}
\definecolor{codestring}{RGB}{5, 150, 105}
\definecolor{codekeyword}{RGB}{99, 102, 241}
\definecolor{lightblue}{RGB}{239, 246, 255}
\definecolor{lightgreen}{RGB}{236, 253, 245}
\definecolor{lightred}{RGB}{254, 242, 242}
\definecolor{jsorange}{RGB}{180, 83, 9}

% ---- Style JavaScript ----
\lstdefinelanguage{JavaScript}{
  keywords={const, let, var, function, return, if, else, new, true, false, null, require, module, exports, async, await, try, catch, throw},
  keywordstyle=\color{codekeyword}\bfseries,
  ndkeywords={Number, String, Boolean, Array, Object, Promise, Date, Math, JSON, Buffer},
  ndkeywordstyle=\color{jsorange}\bfseries,
  sensitive=true,
  comment=[l]{//},
  morecomment=[s]{/*}{*/},
  commentstyle=\color{codecomment}\itshape,
  stringstyle=\color{codestring},
  morestring=[b]',
  morestring=[b]",
  morestring=[b]`,
}

\lstdefinestyle{js}{
  language=JavaScript,
  backgroundcolor=\color{codebg},
  basicstyle=\ttfamily\footnotesize,
  numbers=left,
  stepnumber=1,
  numberstyle=\tiny\color{codecomment},
  numbersep=8pt,
  frame=single,
  rulecolor=\color{codeborder},
  breaklines=true,
  breakatwhitespace=true,
  tabsize=2,
  showstringspaces=false,
  xleftmargin=15pt,
  xrightmargin=5pt,
  captionpos=b,
}

\lstdefinestyle{json}{
  language=,
  backgroundcolor=\color{codebg},
  basicstyle=\ttfamily\footnotesize,
  stringstyle=\color{codestring},
  frame=single,
  rulecolor=\color{codeborder},
  breaklines=true,
  tabsize=2,
  showstringspaces=false,
  xleftmargin=15pt,
  captionpos=b,
}

\lstdefinestyle{bash}{
  language=bash,
  basicstyle=\ttfamily\footnotesize\color{white},
  backgroundcolor=\color[RGB]{30,41,59},
  keywordstyle=\color[RGB]{129,230,217},
  stringstyle=\color[RGB]{167,243,208},
  commentstyle=\color[RGB]{107,114,128}\itshape,
  frame=single,
  rulecolor=\color[RGB]{51,65,85},
  breaklines=true,
  tabsize=2,
  showstringspaces=false,
  xleftmargin=15pt,
  captionpos=b,
}

% ---- En-tete et pied de page ----
\pagestyle{fancy}
\fancyhf{}
\fancyhead[L]{\small\textcolor{primary}{\textbf{Devoir 304 --- API bancaire JavaScript}}}
\fancyhead[R]{\small\textcolor{codecomment}{AMAYEN BOUSSIOM THIERRY MARTIN}}
\fancyfoot[C]{\small\textcolor{codecomment}{\thepage}}
\renewcommand{\headrulewidth}{0.5pt}
\renewcommand{\footrulewidth}{0.5pt}

% ---- Titres ----
\titleformat{\section}{\Large\bfseries\color{primary}}{}{0em}{\thesection\quad}[\titlerule]
\titleformat{\subsection}{\large\bfseries\color{secondary}}{}{0em}{\thesubsection\quad}
\titleformat{\subsubsection}{\normalsize\bfseries\color{secondary}}{}{0em}{\thesubsubsection\quad}

% ---- Boites colorees ----
\newmdenv[backgroundcolor=lightblue,linecolor=primary,linewidth=2pt,
  leftline=true,rightline=false,topline=false,bottomline=false,
  innerleftmargin=10pt,innerrightmargin=10pt,innertopmargin=8pt,innerbottommargin=8pt]{infobox}

\newmdenv[backgroundcolor=lightgreen,linecolor=success,linewidth=2pt,
  leftline=true,rightline=false,topline=false,bottomline=false,
  innerleftmargin=10pt,innerrightmargin=10pt,innertopmargin=8pt,innerbottommargin=8pt]{successbox}

\newmdenv[backgroundcolor=lightred,linecolor=danger,linewidth=2pt,
  leftline=true,rightline=false,topline=false,bottomline=false,
  innerleftmargin=10pt,innerrightmargin=10pt,innertopmargin=8pt,innerbottommargin=8pt]{dangerbox}

\hypersetup{colorlinks=true,linkcolor=primary,urlcolor=primary,
  pdftitle={Rapport Implementation - API Bancaire JavaScript},
  pdfauthor={AMAYEN BOUSSIOM THIERRY MARTIN}}

\setlength{\parskip}{6pt}
\setlength{\parindent}{0pt}

% =========================================================
\begin{document}
% =========================================================

% ---- Page de garde ----
\begin{titlepage}
\centering
\vspace*{1cm}
{\Large\bfseries UNIVERSITE DE YAOUNDE I\\[0.3cm]}
{\large Faculte des Sciences\\[0.1cm]}
{\large Departement d'Informatique\\[2cm]}

\noindent\rule{\linewidth}{2pt}\\[0.5cm]
{\Huge\bfseries\color{primary} Rapport d'Implementation}\\[0.3cm]
{\LARGE\bfseries API de Transaction Bancaire}\\[0.3cm]
{\large Devoir 304 --- Conception d'API REST en JavaScript}\\[0.3cm]
\noindent\rule{\linewidth}{2pt}\\[2cm]

\begin{tabular}{ll}
\textbf{Etudiant :} & AMAYEN BOUSSIOM THIERRY MARTIN \\[0.3cm]
\textbf{Matricule :} & 25I2177 \\[0.3cm]
\textbf{Niveau :} & Licence 3 ICT \\[0.3cm]
\textbf{Etablissement :} & Universite de Yaounde I \\[0.3cm]
\textbf{Depot GitHub :} & \url{https://github.com/AmayenICTL3/TPICT304} \\[0.3cm]
\textbf{Date :} & Avril 2026 \\
\end{tabular}

\vfill
{\large\textit{<<~Conception et implementation d'une API REST en JavaScript\\
pour la gestion de comptes bancaires et de transactions.~>>}}
\end{titlepage}

% ---- Table des matieres ----
\tableofcontents
\newpage

% =========================================================
\section{Introduction et contexte}
% =========================================================

\subsection{Contexte du projet}

Ce rapport decrit de facon complete l'implementation d'une \textbf{API REST de transaction bancaire en JavaScript} realisee dans le cadre du Devoir 304 de Licence 3 ICT.

L'API permet a des clients de gerer leurs comptes bancaires : creer un compte, consulter son solde, effectuer des depots et des retraits, et visualiser l'historique de leurs operations.

\subsection{Probleme resolu}

Les etablissements financiers ont besoin de solutions numeriques pour automatiser la gestion des operations bancaires. Cette API constitue le back-end d'un tel systeme, consommable par une application web, mobile ou desktop.

\subsection{Choix technologiques}

\begin{center}
\begin{tabular}{lll}
\toprule
\textbf{Composant} & \textbf{Technologie} & \textbf{Raison du choix} \\
\midrule
Langage & JavaScript (Node.js) & Ubiquite, asynchrone natif, ecosystem NPM \\
Serveur HTTP & Module natif \texttt{http} & Aucune dependance externe, pedagogique \\
Documentation & Swagger UI (OpenAPI 3.0) & Standard industriel, interface interactive \\
Persistance & Memoire vive (tableau) & Simple, sans configuration, adapte au devoir \\
Identifiants & Module natif \texttt{crypto} & UUID v4 generes sans bibliotheque externe \\
\bottomrule
\end{tabular}
\end{center}

\begin{infobox}
\textbf{Pourquoi JavaScript pur sans framework ?} L'objectif pedagogique est de comprendre le fonctionnement d'un serveur HTTP de l'interieur, sans la magie d'un framework comme Express ou Fastify. Le module \texttt{http} de Node.js est utilise directement.
\end{infobox}

% =========================================================
\section{Architecture du projet}
% =========================================================

\subsection{Structure des fichiers}

\begin{lstlisting}[style=bash, caption={Arborescence du projet}]
TPICT304/
|
|-- server.js              # Serveur HTTP : toute la logique de l'API
|-- openapi.json           # Specification OpenAPI 3.0 (Swagger)
|-- package.json           # Metadonnees Node.js
|-- CAHIER_DE_CHARGES.md   # Cahier des charges
|-- CAHIER_DE_CHARGES.pdf  # Version PDF du cahier des charges
|
|-- public/
    |-- swagger-ui/
        |-- swagger-ui-bundle.js  # Interface Swagger UI
        |-- swagger-ui.css
        |-- favicon.svg
\end{lstlisting}

\subsection{Architecture de server.js}

Contrairement a un projet Python multi-fichiers, toute la logique est contenue dans \textbf{un seul fichier} \texttt{server.js}. Il est organise en blocs fonctionnels distincts :

\begin{center}
\begin{tabular}{clp{8cm}}
\toprule
\textbf{Bloc} & \textbf{Contenu} & \textbf{Role} \\
\midrule
1 & Utilitaires & \texttt{json()}, \texttt{sendFile()}, \texttt{readJsonBody()}, \texttt{round()} \\
2 & Logique metier & \texttt{generateAccountNumber()}, \texttt{accountSummary()}, \texttt{findAccount()} \\
3 & Interface Swagger & \texttt{docsHtml()} \\
4 & Routeur HTTP & \texttt{http.createServer()} avec gestion de toutes les routes \\
\bottomrule
\end{tabular}
\end{center}

% =========================================================
\section{Specifications fonctionnelles}
% =========================================================

\subsection{Acteurs du systeme}

\textbf{Client bancaire} : l'utilisateur principal. Il peut :
\begin{itemize}[noitemsep]
  \item ouvrir un compte bancaire
  \item consulter les informations de son compte
  \item effectuer des depots et des retraits
  \item consulter l'historique de ses transactions
\end{itemize}

\textbf{Administrateur systeme} : supervise le systeme. Il peut :
\begin{itemize}[noitemsep]
  \item consulter la liste complete des comptes
  \item suivre les transactions
\end{itemize}

\subsection{Les 5 fonctions du systeme}

\begin{center}
\begin{tabular}{clll}
\toprule
\textbf{\#} & \textbf{Fonction} & \textbf{Methode HTTP} & \textbf{URL} \\
\midrule
F1 & Creer un compte & POST & \texttt{/accounts} \\
F2 & Lister les comptes & GET & \texttt{/accounts} \\
F3 & Consulter un compte & GET & \texttt{/accounts/\{id\}} \\
F4 & Effectuer un depot & POST & \texttt{/accounts/\{id\}/deposit} \\
F5 & Effectuer un retrait & POST & \texttt{/accounts/\{id\}/withdraw} \\
F+ & Historique des transactions & GET & \texttt{/accounts/\{id\}/transactions} \\
\bottomrule
\end{tabular}
\end{center}

\subsection{Regles de gestion}

\begin{enumerate}[noitemsep]
  \item Chaque compte possede un identifiant unique (UUID)
  \item Chaque compte possede un numero unique au format \texttt{ACC-AAAAMMJJ-XXXXXXXX}
  \item Le solde initial doit etre superieur ou egal a zero
  \item Un depot doit avoir un montant strictement positif
  \item Un retrait doit avoir un montant strictement positif
  \item Un retrait est refuse si le solde est insuffisant
  \item Chaque transaction modifie immediatement le solde
  \item Chaque transaction est conservee dans l'historique du compte
\end{enumerate}

% =========================================================
\section{Specifications non fonctionnelles}
% =========================================================

\begin{center}
\begin{tabular}{lp{5cm}p{6cm}}
\toprule
\textbf{Critere} & \textbf{Exigence} & \textbf{Comment c'est respecte} \\
\midrule
Simplicite & Code lisible et pedagogique & Un seul fichier, fonctions courtes et bien nommees \\
Performance & Reponses rapides & Donnees en memoire, pas d'I/O disque \\
Fiabilite & Donnees correctes en execution & Tableau JavaScript synchronise, pas de race condition \\
Portabilite & Fonctionne sans dependances & Module \texttt{http} natif, seul \texttt{crypto} utilise \\
Maintenabilite & Code structure clairement & Blocs distincts : utilitaires, metier, routeur \\
Disponibilite & API accessible & Fonctionne tant que le processus Node.js est actif \\
\bottomrule
\end{tabular}
\end{center}

% =========================================================
\section{Implementation detaillee --- server.js}
% =========================================================

\subsection{Initialisation et configuration}

\begin{lstlisting}[style=js, caption={server.js --- Imports et configuration}]
const http = require("http");       // serveur HTTP natif
const fs = require("fs");           // lecture de fichiers (Swagger)
const path = require("path");       // gestion des chemins
const { randomUUID } = require("crypto");  // generation UUID v4

const HOST = "0.0.0.0";            // ecoute sur toutes les interfaces
const PORT = Number(process.env.PORT || process.argv[2] || 8003);
// PORT : variable d'environnement > argument CLI > defaut 8003

const accounts = [];   // tableau en memoire : stocke tous les comptes
\end{lstlisting}

\begin{infobox}
\textbf{Persistance en memoire :} Les donnees sont stockees dans le tableau \texttt{accounts} pendant l'execution du serveur. Elles sont perdues a l'arret. C'est un choix delibere pour ce devoir academique (pas de base de donnees requise).
\end{infobox}

\subsection{Fonctions utilitaires}

\subsubsection{Envoyer une reponse JSON}

\begin{lstlisting}[style=js, caption={server.js --- Fonction json()}]
function json(response, statusCode, payload) {
  const body = JSON.stringify(payload);   // serialise l'objet en chaine JSON
  response.writeHead(statusCode, {
    "Content-Type": "application/json; charset=utf-8",
    "Content-Length": Buffer.byteLength(body),  // taille en octets
  });
  response.end(body);  // envoie la reponse et ferme la connexion
}
\end{lstlisting}

Cette fonction est appelee partout dans le routeur pour retourner des reponses JSON avec le bon code HTTP et les bons en-tetes.

\subsubsection{Lire le corps JSON d'une requete}

\begin{lstlisting}[style=js, caption={server.js --- Fonction readJsonBody()}]
function readJsonBody(request) {
  return new Promise((resolve, reject) => {
    let raw = "";

    // Node.js recoit le corps en morceaux (chunks) de facon asynchrone
    request.on("data", (chunk) => {
      raw += chunk;   // concatene les morceaux
    });

    request.on("end", () => {
      if (!raw) {
        resolve({});   // corps vide -> objet vide
        return;
      }
      try {
        resolve(JSON.parse(raw));   // parse le JSON
      } catch (error) {
        reject(error);   // JSON invalide -> erreur
      }
    });

    request.on("error", reject);
  });
}
\end{lstlisting}

\begin{infobox}
\textbf{Pourquoi une Promise ?} HTTP en Node.js est asynchrone par nature : le corps d'une requete arrive en morceaux successifs. La Promise permet d'attendre que tous les morceaux soient recus avant de parser le JSON, en utilisant \texttt{await}.
\end{infobox}

\subsubsection{Fonctions metier}

\begin{lstlisting}[style=js, caption={server.js --- Fonctions metier}]
function round(value) {
  // Arrondit a 2 decimales pour eviter les erreurs flottantes
  // ex: 0.1 + 0.2 = 0.30000000000000004 en JS
  return Math.round(Number(value) * 100) / 100;
}

function generateAccountNumber() {
  // Genere: ACC-20260422-A1B2C3D4
  const datePart = new Date().toISOString().slice(0, 10).replace(/-/g, "");
  const randomPart = randomUUID().replace(/-/g, "").slice(0, 8).toUpperCase();
  return `ACC-${datePart}-${randomPart}`;
}

function accountSummary(account) {
  // Retourne le compte SANS les transactions (pour GET /accounts)
  return {
    id: account.id,
    account_number: account.account_number,
    full_name: account.full_name,
    phone_number: account.phone_number,
    email: account.email,
    balance: account.balance,
    created_at: account.created_at,
  };
}

function accountDetails(account) {
  // Retourne le compte AVEC les transactions (pour GET /accounts/:id)
  return {
    ...accountSummary(account),   // spread: copie tous les champs du resume
    transactions: account.transactions,
  };
}

function findAccount(accountId) {
  // Cherche un compte par son UUID dans le tableau en memoire
  return accounts.find((account) => account.id === accountId);
  // Retourne undefined si non trouve
}
\end{lstlisting}

\subsection{Le routeur HTTP}

Le serveur utilise un unique gestionnaire de requetes qui analyse la methode HTTP et le chemin URL pour dispatcher vers la bonne logique.

\begin{lstlisting}[style=js, caption={server.js --- Squelette du routeur}]
const server = http.createServer(async (request, response) => {
  const url = new URL(request.url, `http://${HOST}:${PORT}`);
  const pathname = url.pathname;
  // request.method = "GET", "POST", etc.
  // pathname = "/accounts", "/accounts/uuid", etc.

  // === Routes statiques ===
  if (request.method === "GET" && pathname === "/") { ... }
  if (request.method === "GET" && pathname === "/docs") { ... }
  if (request.method === "GET" && pathname === "/openapi.json") { ... }
  if (request.method === "GET" && pathname.startsWith("/public/")) { ... }

  // === Routes dynamiques /accounts ===
  if (pathname === "/accounts" && request.method === "GET") { ... }
  if (pathname === "/accounts" && request.method === "POST") { ... }

  // === Routes avec parametre /accounts/:id ===
  const accountMatch = pathname.match(/^\/accounts\/([^/]+)$/);
  // accountMatch[1] = l'UUID extrait de l'URL

  // === Routes avec action /accounts/:id/(deposit|withdraw|transactions) ===
  const txMatch = pathname.match(/^\/accounts\/([^/]+)\/(deposit|withdraw|transactions)$/);
  // txMatch[1] = UUID, txMatch[2] = "deposit" ou "withdraw" ou "transactions"
});
\end{lstlisting}

\begin{infobox}
\textbf{Expressions regulieres pour le routage :} En l'absence de framework, les routes dynamiques sont identifiees par des expressions regulieres. Le groupe de capture \texttt{([{\char`\^}/]+)} extrait l'UUID, et \texttt{(deposit|withdraw|transactions)} identifie l'action.
\end{infobox}

\subsection{Fonction 1 --- Creer un compte (POST /accounts)}

\begin{lstlisting}[style=js, caption={server.js --- POST /accounts}]
if (pathname === "/accounts" && request.method === "POST") {
  try {
    const payload = await readJsonBody(request);  // lit le corps JSON

    // Validation 1 : champs obligatoires
    if (!payload.full_name || !payload.phone_number) {
      json(response, 400, { detail: "full_name et phone_number sont obligatoires." });
      return;
    }

    // Validation 2 : solde initial non negatif
    const initialBalance = round(payload.initial_balance || 0);
    if (initialBalance < 0) {
      json(response, 400, { detail: "Le solde initial doit etre positif ou nul." });
      return;
    }

    // Construction de l'objet compte
    const account = {
      id: randomUUID(),                    // UUID v4 unique
      account_number: generateAccountNumber(),
      full_name: payload.full_name,
      phone_number: payload.phone_number,
      email: payload.email || null,        // null si absent
      balance: initialBalance,
      created_at: new Date().toISOString(),
      transactions: [],                    // tableau vide au depart
    };

    // Si solde initial > 0, on cree une transaction "depot initial"
    if (initialBalance > 0) {
      account.transactions.push({
        transaction_id: randomUUID(),
        transaction_type: "deposit",
        amount: initialBalance,
        description: "Solde initial",
        balance_after: initialBalance,
        created_at: new Date().toISOString(),
      });
    }

    accounts.push(account);               // ajoute au tableau memoire
    json(response, 201, accountSummary(account));  // HTTP 201 Created
  } catch {
    json(response, 400, { detail: "Corps JSON invalide." });
  }
  return;
}
\end{lstlisting}

\textbf{Etapes de traitement :}
\begin{enumerate}[noitemsep]
  \item Lecture asynchrone du corps JSON via \texttt{readJsonBody}
  \item Validation des champs obligatoires (\texttt{full\_name}, \texttt{phone\_number})
  \item Validation du solde initial ($\geq 0$)
  \item Construction de l'objet compte avec UUID et numero generes
  \item Si solde > 0 : ajout d'une transaction initiale dans l'historique
  \item Ajout dans le tableau \texttt{accounts}
  \item Reponse HTTP 201 avec le resume du compte (sans les transactions)
\end{enumerate}

\subsection{Fonction 2 --- Lister les comptes (GET /accounts)}

\begin{lstlisting}[style=js, caption={server.js --- GET /accounts}]
if (pathname === "/accounts" && request.method === "GET") {
  // accounts.map(accountSummary) : transforme chaque compte
  // en version "sans transactions" pour alleger la reponse
  json(response, 200, accounts.map(accountSummary));
  return;
}
\end{lstlisting}

Simple mais efficace : \texttt{Array.map()} applique \texttt{accountSummary} a chaque element du tableau et retourne un nouveau tableau avec uniquement les champs pertinents.

\subsection{Fonction 3 --- Consulter un compte (GET /accounts/\{id\})}

\begin{lstlisting}[style=js, caption={server.js --- GET /accounts/\{account\_id\}}]
const accountMatch = pathname.match(/^\/accounts\/([^/]+)$/);
if (accountMatch && request.method === "GET") {
  const account = findAccount(accountMatch[1]);  // recherche par UUID
  if (!account) {
    json(response, 404, { detail: "Compte introuvable." });
    return;
  }
  json(response, 200, accountDetails(account));  // avec transactions
  return;
}
\end{lstlisting}

\texttt{accountMatch[1]} contient l'UUID extrait de l'URL par l'expression reguliere. Si \texttt{findAccount} retourne \texttt{undefined}, on repond HTTP 404.

\subsection{Fonctions 4 et 5 --- Depot et Retrait}

\begin{lstlisting}[style=js, caption={server.js --- POST /accounts/\{id\}/(deposit|withdraw)}]
const txMatch = pathname.match(
  /^\/accounts\/([^/]+)\/(deposit|withdraw|transactions)$/
);
if (txMatch) {
  const account = findAccount(txMatch[1]);   // recherche le compte
  if (!account) {
    json(response, 404, { detail: "Compte introuvable." });
    return;
  }

  const action = txMatch[2];  // "deposit", "withdraw" ou "transactions"

  if ((action === "deposit" || action === "withdraw") && request.method === "POST") {
    try {
      const payload = await readJsonBody(request);
      const amount = round(payload.amount);

      // Validation : montant strictement positif
      if (!Number.isFinite(amount) || amount <= 0) {
        json(response, 400, { detail: "Le montant doit etre strictement positif." });
        return;
      }

      // Validation specifique au retrait : solde suffisant
      if (action === "withdraw" && amount > account.balance) {
        json(response, 400, { detail: "Solde insuffisant pour effectuer ce retrait." });
        return;
      }

      // Mise a jour du solde
      account.balance = action === "deposit"
        ? round(account.balance + amount)   // depot : on ajoute
        : round(account.balance - amount);  // retrait : on soustrait

      // Enregistrement de la transaction dans l'historique
      account.transactions.push({
        transaction_id: randomUUID(),
        transaction_type: action,
        amount,
        description: payload.description || null,
        balance_after: account.balance,
        created_at: new Date().toISOString(),
      });

      json(response, 200, accountDetails(account));  // compte mis a jour
    } catch {
      json(response, 400, { detail: "Corps JSON invalide." });
    }
    return;
  }
}
\end{lstlisting}

\begin{infobox}
\textbf{Modification directe de l'objet en memoire :} Comme \texttt{account} est une reference vers l'objet dans le tableau \texttt{accounts}, modifier \texttt{account.balance} et \texttt{account.transactions} modifie directement la donnee en memoire. Pas besoin de "sauvegarder" explicitement.
\end{infobox}

% =========================================================
\section{Cas de test}
% =========================================================

\subsection{Tableau complet des 23 cas de test}

\begin{longtable}{clp{3.5cm}p{3cm}c}
\toprule
\textbf{ID} & \textbf{Fonction} & \textbf{Description} & \textbf{Resultat attendu} & \textbf{Code HTTP} \\
\midrule
\endhead
CT-01 & Creer compte & Tous les champs fournis & Compte cree avec ID et numero & 201 \\
CT-02 & Creer compte & Sans email ni solde initial & Compte cree, balance = 0 & 201 \\
CT-03 & Creer compte & \texttt{full\_name} absent & Erreur champ obligatoire & 400 \\
CT-04 & Creer compte & \texttt{phone\_number} absent & Erreur champ obligatoire & 400 \\
CT-05 & Creer compte & Solde initial negatif & Erreur solde invalide & 400 \\
\midrule
CT-06 & Lister comptes & Aucun compte cree & Tableau vide \texttt{[]} & 200 \\
CT-07 & Lister comptes & Comptes existants & Liste des comptes & 200 \\
\midrule
CT-08 & Consulter compte & ID valide & Details + transactions & 200 \\
CT-09 & Consulter compte & ID invalide & Compte introuvable & 404 \\
\midrule
CT-10 & Depot & Montant valide & Solde augmente & 200 \\
CT-11 & Depot & Avec description & Libelle visible & 200 \\
CT-12 & Depot & Montant = 0 & Erreur validation & 400 \\
CT-13 & Depot & Montant negatif & Erreur validation & 400 \\
CT-14 & Depot & Compte inexistant & Compte introuvable & 404 \\
\midrule
CT-15 & Retrait & Solde suffisant & Solde diminue & 200 \\
CT-16 & Retrait & Montant = solde & Solde = 0 & 200 \\
CT-17 & Retrait & Solde insuffisant & Solde insuffisant & 400 \\
CT-18 & Retrait & Montant = 0 & Erreur validation & 400 \\
CT-19 & Retrait & Montant negatif & Erreur validation & 400 \\
CT-20 & Retrait & Compte inexistant & Compte introuvable & 404 \\
\midrule
CT-21 & Transactions & Compte sans operation & Tableau vide & 200 \\
CT-22 & Transactions & Avec operations & Liste chronologique & 200 \\
CT-23 & Transactions & Compte inexistant & Compte introuvable & 404 \\
\bottomrule
\end{longtable}

\subsection{Test manuel detaille --- Scenario complet}

Les tests manuels se font dans \textbf{Swagger UI} a l'adresse \texttt{http://localhost:8003/docs}.

\subsubsection{Etape 1 --- Verifier la liste vide (CT-06)}

\begin{lstlisting}[style=bash, caption={Requete GET /accounts}]
GET http://localhost:8003/accounts
\end{lstlisting}

\begin{successbox}
\textbf{Resultat attendu :} HTTP 200 avec le corps \texttt{[]}\\
Confirme que le serveur repond et que le tableau est vide au demarrage.
\end{successbox}

\subsubsection{Etape 2 --- Creer un compte (CT-01)}

Dans Swagger, ouvrir \textbf{POST /accounts}, cliquer \textbf{Try it out}, selectionner l'exemple \textbf{CT-01} dans le menu deroulant \textit{Examples}.

\begin{lstlisting}[style=json, caption={Corps de la requete}]
{
  "full_name": "AMAYEN BOUSSIOM THIERRY MARTIN",
  "phone_number": "699001122",
  "email": "thierry@example.com",
  "initial_balance": 10000
}
\end{lstlisting}

\begin{lstlisting}[style=json, caption={Reponse attendue --- HTTP 201}]
{
  "id": "a3f1b2c4-d5e6-7890-abcd-ef1234567890",
  "account_number": "ACC-20260422-A1B2C3D4",
  "full_name": "AMAYEN BOUSSIOM THIERRY MARTIN",
  "phone_number": "699001122",
  "email": "thierry@example.com",
  "balance": 10000,
  "created_at": "2026-04-22T10:00:00.000Z"
}
\end{lstlisting}

\begin{infobox}
\textbf{A noter :} Copier l'\texttt{id} retourne pour les etapes suivantes.
\end{infobox}

\subsubsection{Etape 3 --- Tester les cas d'erreur (CT-03, CT-04, CT-05)}

Selectionner l'exemple \textbf{CT-05} dans le menu Examples :
\begin{lstlisting}[style=json, caption={Test CT-05 : solde negatif}]
{
  "full_name": "TEST NEGATIF",
  "phone_number": "600000001",
  "initial_balance": -500
}
\end{lstlisting}

\begin{dangerbox}
\textbf{Resultat attendu :} HTTP 400\\
Corps : \texttt{\{"detail": "Le solde initial doit etre positif ou nul."\}}
\end{dangerbox}

\subsubsection{Etape 4 --- Effectuer un depot (CT-10)}

\begin{lstlisting}[style=bash, caption={Requete POST /accounts/\{id\}/deposit}]
POST http://localhost:8003/accounts/a3f1b2c4-.../deposit
\end{lstlisting}

Selectionner l'exemple \textbf{CT-10} :
\begin{lstlisting}[style=json, caption={Corps}]
{ "amount": 5000 }
\end{lstlisting}

\begin{successbox}
\textbf{Resultat attendu :} HTTP 200, balance = 15000 (10000 + 5000).\\
2 transactions visibles : solde initial + depot.
\end{successbox}

\subsubsection{Etape 5 --- Effectuer un retrait valide (CT-15)}

Selectionner l'exemple \textbf{CT-15} :
\begin{lstlisting}[style=json, caption={Corps}]
{ "amount": 2000, "description": "Retrait client" }
\end{lstlisting}

\begin{successbox}
\textbf{Resultat attendu :} HTTP 200, balance = 13000 (15000 - 2000).
\end{successbox}

\subsubsection{Etape 6 --- Tester le solde insuffisant (CT-17)}

Selectionner l'exemple \textbf{CT-17} :
\begin{lstlisting}[style=json, caption={Corps}]
{ "amount": 50000, "description": "Test solde insuffisant" }
\end{lstlisting}

\begin{dangerbox}
\textbf{Resultat attendu :} HTTP 400\\
Corps : \texttt{\{"detail": "Solde insuffisant pour effectuer ce retrait."\}}
\end{dangerbox}

\subsubsection{Etape 7 --- Consulter l'historique (CT-22)}

\begin{lstlisting}[style=bash, caption={Requete GET /accounts/\{id\}/transactions}]
GET http://localhost:8003/accounts/a3f1b2c4-.../transactions
\end{lstlisting}

\begin{successbox}
\textbf{Resultat attendu :} HTTP 200 avec 3 transactions dans l'ordre chronologique :\\
1. Depot initial (10 000) --- balance 10 000\\
2. Depot (5 000) --- balance 15 000\\
3. Retrait (2 000) --- balance 13 000
\end{successbox}

\subsubsection{Etape 8 --- Tester un ID inexistant (CT-09)}

\begin{lstlisting}[style=bash, caption={Test CT-09}]
GET http://localhost:8003/accounts/00000000-0000-0000-0000-000000000000
\end{lstlisting}

\begin{dangerbox}
\textbf{Resultat attendu :} HTTP 404\\
Corps : \texttt{\{"detail": "Compte introuvable."\}}
\end{dangerbox}

% =========================================================
\section{Documentation Swagger}
% =========================================================

\subsection{Fonctionnement}

L'interface Swagger UI est accessible a \texttt{http://localhost:8003/docs}. Elle est servie par \texttt{server.js} depuis le dossier \texttt{public/swagger-ui/} et se nourrit du fichier \texttt{openapi.json} pour afficher la documentation interactive.

\begin{lstlisting}[style=js, caption={server.js --- Generation de la page Swagger}]
function docsHtml() {
  return `<!DOCTYPE html>
<html>
<head>
  <link rel="stylesheet" href="/public/swagger-ui/swagger-ui.css">
  <title>API JavaScript - Swagger UI</title>
</head>
<body>
  <div id="swagger-ui"></div>
  <script src="/public/swagger-ui/swagger-ui-bundle.js"></script>
  <script>
    SwaggerUIBundle({
      url: '/openapi.json',          // charge le fichier de spec
      dom_id: '#swagger-ui',
      deepLinking: true,
      docExpansion: 'list',
      displayRequestDuration: true,  // affiche le temps de reponse
      filter: true,
    });
  </script>
</body>
</html>`;
}
\end{lstlisting}

\subsection{Exemples interactifs dans Swagger}

Le fichier \texttt{openapi.json} contient des exemples nommes pour chaque endpoint :

\begin{center}
\begin{tabular}{ll}
\toprule
\textbf{Endpoint} & \textbf{Exemples disponibles dans Swagger} \\
\midrule
\texttt{POST /accounts} & CT-01, CT-02, CT-03, CT-04, CT-05 \\
\texttt{POST /\{id\}/deposit} & CT-10, CT-11, CT-12, CT-13 \\
\texttt{POST /\{id\}/withdraw} & CT-15, CT-16, CT-17, CT-18, CT-19 \\
\bottomrule
\end{tabular}
\end{center}

\subsection{Comment utiliser Swagger pour les tests manuels}

\begin{enumerate}
  \item Lancer le serveur : \texttt{node server.js}
  \item Ouvrir \texttt{http://localhost:8003/docs} dans le navigateur
  \item Cliquer sur un endpoint (ex : \textit{POST /accounts})
  \item Cliquer sur \textbf{Try it out}
  \item Selectionner un cas de test dans le menu deroulant \textbf{Examples}
  \item Cliquer sur \textbf{Execute}
  \item Verifier le code HTTP et le corps de la reponse
\end{enumerate}

% =========================================================
\section{Lancement et deploiement}
% =========================================================

\subsection{Lancement en local}

\begin{lstlisting}[style=bash, caption={Lancement du serveur}]
# Lancer le serveur sur le port par defaut (8003)
node server.js

# Ou sur un port personnalise
node server.js 4000

# L'API est disponible sur :
#   http://localhost:8003/        (point d'entree)
#   http://localhost:8003/docs    (Swagger UI)
\end{lstlisting}

\subsection{package.json}

\begin{lstlisting}[style=json, caption={package.json}]
{
  "name": "api-bancaire-javascript",
  "version": "1.0.0",
  "description": "API REST bancaire en JavaScript pur",
  "main": "server.js",
  "scripts": {
    "start": "node server.js"
  }
}
\end{lstlisting}

Aucune dependance NPM : \texttt{npm install} n'est pas necessaire.

% =========================================================
\section{Conclusion}
% =========================================================

Ce projet implemente une \textbf{API REST complete en JavaScript pur} pour la gestion de comptes bancaires, sans aucune dependance externe, respectant toutes les exigences fonctionnelles et non fonctionnelles du cahier des charges.

\subsection*{Points cles de l'implementation}

\begin{itemize}[noitemsep]
  \item \textbf{JavaScript pur :} module \texttt{http} natif, aucune dependance NPM externe
  \item \textbf{Routage par regex :} expressions regulieres pour extraire les parametres d'URL
  \item \textbf{Asynchronisme :} \texttt{async/await} et Promises pour la lecture du corps HTTP
  \item \textbf{Validation manuelle :} toutes les regles de gestion implementees explicitement
  \item \textbf{Documentation interactive :} 23 cas de test executables dans Swagger UI
  \item \textbf{Codes HTTP corrects :} 201 (creation), 200 (succes), 400 (erreur metier), 404 (non trouve)
\end{itemize}

\subsection*{Limites actuelles}

\begin{itemize}[noitemsep]
  \item Pas de persistance : donnees perdues a l'arret du serveur
  \item Pas d'authentification
  \item Pas de virement entre comptes
  \item Pas de separation des roles client/administrateur
\end{itemize}

Ces elements representent des evolutions naturelles pour une version future.

\vfill
\begin{center}
\rule{0.5\linewidth}{0.5pt}\\[0.3cm]
{\small AMAYEN BOUSSIOM THIERRY MARTIN --- Matricule 25I2177 --- Licence 3 ICT\\
Universite de Yaounde I --- Avril 2026}
\end{center}

\end{document}
